% =========================================================
% Section 3 — Sequential Characterizations (Bridge)
%
% Source: volume-iii/analysis/metric-spaces/notes/
% Purpose: Bridge from metric geometry to topological behaviour
%          via sequential characterizations of open and closed sets.
%
% Dependencies:
%   open-closed-sets.tex     — def:open-closed-set (ball-based definition),
%                              def:closure, def:adherent-point
%   \usepackage{tikz}
%   \usetikzlibrary{arrows.meta,positioning}
%   \usepackage{booktabs}
%   \usepackage{enumitem}
% =========================================================

\subsection{Sequential Characterizations of Open and Closed Sets}
\label{sec:sequential-characterizations}

In metric spaces, properties of sets can often be described
in terms of the behaviour of sequences.
In particular, the ball-based definitions of open and closed sets
established in \hyperref[sec:open-closed-definitions]{\S\ref*{sec:open-closed-definitions}}
admit equivalent sequential formulations that link metric geometry
to topological structure.
These are \emph{characterizations}, not new definitions: each statement
below is provably equivalent to the ball-based definition.

% ---------------------------------------------------------
\subsubsection{Closed Sets via Sequences}
% ---------------------------------------------------------

\begin{tcolorbox}[colback=propbox, colframe=propborder, arc=2pt,
  left=6pt, right=6pt, top=4pt, bottom=4pt,
  title={\small\textbf{Sequential Characterization (Closed Sets)}},
  fonttitle=\small\bfseries]
\label{def:closed-sequential}
Let $(X,d)$ be a metric space and let $F \subseteq X$.

\medskip
$F$ is closed (i.e.\ $X \setminus F$ is open in the sense of
\hyperref[def:open-closed-set]{Definition~\ref*{def:open-closed-set}})
if and only if every convergent sequence of points of $F$ converges to a
point of $F$.
\end{tcolorbox}

\paragraph{Compact form.}

\[
F \text{ is closed}
\iff
\forall (x_n)_{n\in\mathbb{N}} \subseteq F,
\Bigl(
  x_n \to x
  \;\Rightarrow\;
  x \in F
\Bigr).
\]

\paragraph{Fully quantified form.}

Expanding the definition of convergence in full:

\[
F \text{ is closed}
\iff
\forall (x_n)_{n\in\mathbb{N}} \subseteq F,\;
\forall x \in X,\;
\left(
  \Bigl[
    \forall \varepsilon>0\;
    \exists N\in\mathbb{N}\;
    \forall n\ge N,\;
    d(x_n,x)<\varepsilon
  \Bigr]
  \Rightarrow
  x\in F
\right).
\]

\begin{remark*}
This formulation captures the idea that closed sets
\emph{contain all of their limit points}.
Limits of sequences from the set cannot ``escape'' the set.
\end{remark*}

% ---------------------------------------------------------
\subsubsection{Open Sets via Sequences}
% ---------------------------------------------------------

\begin{tcolorbox}[colback=propbox, colframe=propborder, arc=2pt,
  left=6pt, right=6pt, top=4pt, bottom=4pt,
  title={\small\textbf{Sequential Characterization (Open Sets)}},
  fonttitle=\small\bfseries]
\label{def:open-sequential}
Let $(X,d)$ be a metric space and let $U \subseteq X$.

\medskip
$U$ is open (i.e.\ satisfies \hyperref[def:open-closed-set]{Definition~\ref*{def:open-closed-set}})
if and only if every sequence that converges to a point of $U$ eventually
lies entirely in $U$.
\end{tcolorbox}

\paragraph{Compact form.}

\[
U \text{ is open}
\iff
\forall x\in U,\;
\forall (x_n)_{n\in\mathbb{N}} \subseteq X,\;
\Bigl(
  x_n \to x
  \;\Rightarrow\;
  \exists N\in\mathbb{N}\;
  \forall n\ge N,\;
  x_n\in U
\Bigr).
\]

\paragraph{Fully quantified form.}

Expanding the convergence condition in full:

\[
U \text{ is open}
\iff
\forall x\in U,\;
\forall (x_n)\subseteq X,\;
\left(
  \Bigl[
    \forall \varepsilon>0\;
    \exists N\in\mathbb{N}\;
    \forall n\ge N,\;
    d(x_n,x)<\varepsilon
  \Bigr]
  \Rightarrow
  \exists N'\in\mathbb{N}\;
  \forall n\ge N',\;
  x_n\in U
\right).
\]

\begin{remark*}
This condition expresses the idea that points of an open set
are \emph{stable under small perturbations}.
If a sequence approaches a point inside the set, the sequence
must eventually remain inside the set.
\end{remark*}

% ---------------------------------------------------------
\subsubsection{Duality Between Open and Closed Sets}
% ---------------------------------------------------------

The sequential characterizations of open and closed sets
are complementary.

\begin{center}
\begin{tabular}{cc}
\toprule
\textbf{Closed Sets} & \textbf{Open Sets} \\
\midrule
Sequences inside the set & Sequences converging to a point \\
keep their limits inside & of the set eventually enter it \\
\bottomrule
\end{tabular}
\end{center}

\noindent Symbolically,

\[
\textbf{Closed:}
\qquad
x_n \in F,\; x_n \to x
\;\Rightarrow\;
x\in F
\]

\[
\textbf{Open:}
\qquad
x_n \to x \in U
\;\Rightarrow\;
x_n \in U
\text{ eventually.}
\]

\begin{remark*}[Conceptual Summary]
These two statements reflect a fundamental duality.
\begin{itemize}
\item Closed sets are \emph{stable under limits}.
\item Open sets are \emph{stable under small perturbations}.
\end{itemize}
Both viewpoints rely on the structure of metric spaces,
where sequences are powerful enough to detect topological properties.
\end{remark*}

\begin{remark*}[Why This Works in Metric Spaces]
\label{rem:first-countable}
The sequential characterizations above rely on the fact that metric
spaces are \emph{first-countable}.
In such spaces, convergence of sequences is sufficient to capture the
topological notions of openness and closedness.
In more general topological spaces, sequences are not always sufficient
to detect these properties.
\end{remark*}

% ---------------------------------------------------------
\subsubsection{Equivalent Characterizations of Closed Sets}
% ---------------------------------------------------------

\begin{theorem}[Equivalent Characterizations of Closed Sets]
\label{thm:closed-equiv}
Let $(X,d)$ be a metric space and let $F \subseteq X$.
The following statements are equivalent.

\begin{enumerate}[label=(\roman*)]
\item $F$ is closed in the sense of
      \hyperref[def:open-closed-set]{Definition~\ref*{def:open-closed-set}}:
      $X \setminus F$ is open.

\item $F$ contains all of its limit points.

\item $F$ is sequentially closed: for every sequence
      $(x_n) \subseteq F$,
      \[
        x_n \to x \quad\Rightarrow\quad x \in F.
      \]
\end{enumerate}
\end{theorem}

\begin{remark*}
These equivalences rely on the structure of metric spaces.
In general topological spaces, sequential closure and topological
closure need not coincide.
\end{remark*}

\begin{remark*}[Interpretation]
The three conditions highlight different perspectives on the same property.
\begin{itemize}
  \item \textbf{Topological.} Closed sets are complements of open sets.
  \item \textbf{Limit-point.} Closed sets contain all of their
        accumulation points.
  \item \textbf{Sequential.} Limits of sequences in the set remain
        in the set.
\end{itemize}
Together these viewpoints provide a powerful toolkit for reasoning
about closed sets in analysis.
\end{remark*}

% ---------------------------------------------------------
% TikZ: Closed-set equivalence triangle
% ---------------------------------------------------------

\begin{center}
\begin{tikzpicture}[
  >=Latex,
  node distance=22mm,
  box/.style={draw, rounded corners=3pt, align=center, inner sep=6pt},
  lab/.style={font=\small},
  eq/.style={<->, thick}
]

\node[box] (closed) {\textbf{Closed}\\[-1pt]\scriptsize $F$ is closed};

\node[box, below left=of closed] (limitpt)
  {\textbf{Limit-point}\\[-1pt]\scriptsize $F$ contains all\\[-1pt]\scriptsize limit points};

\node[box, below right=of closed] (seq)
  {\textbf{Sequential}\\[-1pt]\scriptsize $(x_n)\subseteq F,\ x_n\to x$\\[-1pt]\scriptsize $\Rightarrow\ x\in F$};

\draw[eq] (closed)  -- node[lab, left]  {$\Longleftrightarrow$} (limitpt);
\draw[eq] (closed)  -- node[lab, right] {$\Longleftrightarrow$} (seq);
\draw[eq] (limitpt) -- node[lab, below] {$\Longleftrightarrow$} (seq);

\node[lab, above=4mm of closed]
  {Closed set characterizations (metric spaces)};

\end{tikzpicture}
\end{center}

% ---------------------------------------------------------
\subsubsection{Sequential Duality Proposition}
% ---------------------------------------------------------

\begin{proposition}[Sequential Duality of Open and Closed Sets]
\label{prop:sequential-duality}
Let $(X,d)$ be a metric space, let $U \subseteq X$, and let
$F = X \setminus U$.
Then

\[
x_n \in F,\; x_n \to x
\;\Rightarrow\;
x \in F
\quad\Longleftrightarrow\quad
x_n \to x \in U
\;\Rightarrow\;
x_n \in U \text{ eventually.}
\]
\end{proposition}

\begin{remark*}
The left-hand side expresses the defining sequential property of
\emph{closed sets}: limits of sequences in $F$ remain in $F$.

The right-hand side expresses the defining sequential property of
\emph{open sets}: sequences converging to a point of $U$ must
eventually lie entirely inside $U$.

The biconditional follows by taking complements and applying De Morgan's
laws:
\[
(A \cap B)^c = A^c \cup B^c,
\qquad
(A \cup B)^c = A^c \cap B^c.
\]

The compact summary is

\[
\boxed{\text{Closed = limit stability}}
\qquad
\boxed{\text{Open = neighborhood stability.}}
\]
\end{remark*}

% ---------------------------------------------------------
% TikZ: Open/Closed sequential symmetry diagram
% ---------------------------------------------------------

\begin{center}
\begin{tikzpicture}[
  >=Latex,
  node distance=3.8cm,
  box/.style={draw, rounded corners=3pt, align=center, inner sep=8pt},
  lab/.style={font=\small}
]

\node[box] (closed) {
  \textbf{Closed Sets}\\[4pt]
  $(x_n)\subseteq F$\\
  $x_n \to x$\\
  $\Rightarrow x\in F$
};

\node[box, right of=closed] (open) {
  \textbf{Open Sets}\\[4pt]
  $x_n \to x\in U$\\
  $\Rightarrow$ eventually\\
  $x_n\in U$
};

\draw[<->, thick] (closed) -- node[lab, above] {complements / duality} (open);

\node[lab, below=1.4cm of closed] {limits stay inside};

\node[lab, below=1.4cm of open] {approaching sequences enter};

\node[lab, below=0.4cm of open, xshift=-1.9cm]
  {$\boxed{\text{Closed: stable under limits}}$};

\node[lab, below=0.9cm of open, xshift=-1.9cm]
  {$\boxed{\text{Open: stable under small perturbations}}$};

\end{tikzpicture}
\end{center}